\subsection{A bounded all-order test of real-multiplication propagation}
\label{sec:kt-rm-theta-propagation-limit}

This note records an obstruction to one attempted incidence argument.
It does not produce a cyclic cover, a counterexample to the universal
cover question, or a new bound in A45--A52.

\begin{proposition}[KT-RM-THETA-STABILIZER]
Let $k$ be an algebraically closed field of characteristic two, let
$X/k$ be an ordinary smooth projective genus-two curve, and let
$A=\operatorname{Pic}^0(X)$.  Let $B$ be the actual locally exact
differential line bundle and
\[
 \Theta=\{\mathcal O_X(R)\otimes B^{-1}:R\in X(k)\}\subset A.
                                                               \tag{RT.1}
\]
Suppose $\mathcal O_5\hookrightarrow\operatorname{End}(A)$ is
a unital integral action fixed by the Rosati involution for the
canonical principal polarization.  If $0\ne u\in\mathcal O_5$
and $a\in A(k)$ satisfy
\[
                         u(\Theta)\subset\Theta+a,
                                                               \tag{RT.2}
\]
then $u=1$ or $u=-1$, and $a=0$.
Consequently, for every fixed $u\in\mathcal O_5\setminus\{0,1,-1\}$
and every $a\in A(k)$, the set
\[
                    \{z\in\Theta(k):u z\in\Theta+a\}
                                                               \tag{RT.3}
\]
is finite.  This conclusion concerns all torsion orders at once,
but does not exclude choosing exceptional points in that finite set.
\end{proposition}

\begin{proof}
Since $u\bar u=N_{\mathbf Q(\sqrt5)/\mathbf Q}(u)$ is a nonzero
integer, $u$ is an isogeny.  Thus (RT.2) is equality of reduced
curves, and restriction gives a finite surjective morphism of
genus-two curves of some degree $d\ge1$.  Write $D=\deg_A(u)$
for the degree of the isogeny, which need not equal $d$.

Let $h=[\Theta]\in\operatorname{NS}(A)\otimes\mathbf Q$.
Translation acts trivially on this group, and proper pushforward
therefore gives $u_*h=d h$.  The finite-flat identity
$u_*u^*=D$ shows that $u_*$ is invertible on this finite-dimensional
space.  Hence
\[
                       u^*h=(D/d)h.                          \tag{RT.4}
\]
The polarization map $[L]\mapsto\lambda_\Theta^{-1}\phi_L$
is injective on $\operatorname{NS}(A)\otimes\mathbf Q$.
The pullback identity
$\phi_{u^*L}=u^\vee\phi_Lu$ sends $u^*h$ to $u^\dagger u=u^2$.
It follows that
\[
                   u^2=c\in\mathbf Q_{>0},\qquad c=D/d.
                                                               \tag{RT.5}
\]
Since $u$ is an algebraic integer, $c$ is a positive integer.
Degrees on a surface give
$D^2=\deg(u^2)=\deg([c])=c^4$, so $D=c^2$ and $d=c$.
In particular $u^2=d$.

Factor the curve morphism as a power of relative Frobenius
followed by a separable map.  The Frobenius twist still has
genus two.  Riemann--Hurwitz for the separable factor gives
$2\ge2\deg(f_{\mathrm{sep}})$, so that factor has degree one
and is an isomorphism.  Therefore $d=2^r$ for some $r\ge0$,
and the original curve map is radicial.  The equation
$u^2=2^r$ in $\mathbf Q(\sqrt5)$ forces $r=2b$ and
$u=\pm2^b$: for odd $r$ it would put $\sqrt2$ in
$\mathbf Q(\sqrt5)$.

An ordinary genus-two curve in characteristic two has three
distinct Weierstrass points $W_0,W_1,W_2$.  Their classes
$z_i=\mathcal O_X(W_i)\otimes B^{-1}$ are three distinct
points of $\Theta\cap A[2]$, because $2W_i\sim K_X\sim2B$.
They are nonzero because $B$ is ineffective in the ordinary
case.  These are the carried A35, (T2F.3)--(T2F.4), facts for
the actual Raynaud translate.  If $b>0$, $u=\pm[2^b]$ sends
all three $z_i$ to zero, contradicting injectivity of a
radicial curve map on geometric points.  Thus $b=0$ and
$u=\pm1$.

The identity $2B\sim K_X$ makes $\Theta$ symmetric under
$[-1]$, by the hyperelliptic involution.  Hence (RT.2) now says
$\Theta+a=\Theta$.  Such an $a$ lies in the kernel of
$\phi_{\mathcal O_A(\Theta)}$, which is trivial because this
polarization is principal.  Thus $a=0$.

Finally, if (RT.3) were infinite, the closed subset
$\Theta\cap u^{-1}(\Theta+a)$ would be all of the irreducible
projective curve $\Theta$.  This would give (RT.2), already
excluded for the indicated $u$.
\end{proof}

\paragraph{The attempted arithmetic implication and its exact failure.}
Write $T^2+T-1=0$ for the chosen generator of $\mathcal O_5$.
The proposed propagation was to start with odd torsion on theta,
use powers of $T$ to produce more incidences, and combine these
with Frobenius to meet all primitive doubling orbits of one
cyclic subgroup.  Proposition KT-RM-THETA-STABILIZER rules out
the required global preservation of theta by every nontrivial
power of $T$, even allowing a translation.  A finite collection
of such endomorphisms has only finitely many exceptional
incidences of the indicated form.  An incidence at one specially
selected torsion point is still possible and has not been excluded.

Arithmetic Frobenius does preserve the actual theta curve over
its field of definition, but it need not preserve the cyclic
subgroup generated by a chosen theta point.  Even if a chosen
point $z$ has exact order $n$ and satisfies $\Phi z=vz$ for
$v\in(\mathbf Z/n)^\times$, the elementary sufficient condition
$\langle2,-1,v\rangle=(\mathbf Z/n)^\times$ still requires
that particular theta point.  The real-multiplication relation
does not supply it.  No existence assertion is inferred from
this conditional observation.

The converse attempt to avoid every odd theta order on one
finite-field curve is already excluded by the carried A42,
(PE.14)--(PE.15): every such theta curve has infinitely many
nonzero odd-order points.  Their full orders and their distribution
inside one cyclic subgroup remain uncontrolled.  The geometric
generic avoidance of A21/A37 cannot be specialized simultaneously
at all odd levels by the finite-level argument.  In particular,
the construction imposing $Q\equiv1\pmod m$ for every prescribed
finite $m$ supplies varying finite fields; one fixed finite
$Q$ cannot satisfy these congruences for all odd $m$.

These limitations coexist with the carried ordinary
$\mathcal O_5$ example A13/A33 that has a degree-$4159$
primitive two-rank-zero cover.  Thus they imply neither failure
of all odd cyclic covers nor their existence on an arbitrary
ordinary $\mathcal O_5$ curve.  This attempt is classified
\emph{nonmovement} for the universal Route~2 gate.

\paragraph{Exact inputs and proof-reading scope.}
For the elementary abelian-variety inputs see Milne,
\emph{Abelian Varieties}, version 2 January 2022:
Proposition~8.1 and Theorem~8.2, pp.~12--14;
Corollary~12.8, p.~22; Remark~12.13, p.~24;
Lemma~16.2 and its proof, p.~30; and Proposition~17.2 and
its proof, p.~35.  The finite-flat push--pull identity is
Stacks, Section~42.15, Lemma~42.15.2 (tag~\texttt{02RF});
proper pushforward and flat pullback preserve algebraic
equivalence by their compatibility with families in the same
section.  The curve factorization and genus preservation are
Stacks, Section~53.13, Proposition~53.13.7 and
Lemma~53.13.5 (tag~\texttt{0CCV}); Riemann--Hurwitz is
Section~53.12, Lemma~53.12.2 (tag~\texttt{0C1B}).
The statements and the displayed proofs of these local inputs
were read; foundational results cited inside them are imported.
The actual theta-characteristic and two-torsion facts are the
accepted A35 proof, and the odd-torsion abundance is accepted
A42, not newly derived here.  Retrieval metadata and exact
source hashes are retained with this note.
